\section{Monte Carlo evidence for forecast-optimal smoothness}
\label{sec:empirical}

The principal empirical question is whether choosing normalized PLS
smoothness by chronological, horizon-matched forecast MSE is useful
relative to established tuning criteria. The results below use the frozen pooled-selector
experiment. Additional time-adaptive experiments,
with distinct datasets and loss scales, are
reported separately in the appendices.

\subsection{Horizon-matched tuning and forecast error}

The pooled-selector simulation (Checkpoint 03) comprises 3,000 synthetic
scenarios, formed from five latent trends, three noise mechanisms, two noise
scales, and 100 random seeds. Each scenario has six chronological forecast
origins evaluated at horizons $h\in\{1,3,6,12\}$, giving 72,000
outer-origin--horizon evaluations. The estimator fixes $d=2$ and $L=60$.

\Cref{tab:cp03} reports the geometric root-mean-square forecast error
(RMSFE) ratio of horizon-matched pooled forecast-CV to two feasible selectors.
At $h=1$, horizon-matched and one-step forecast-CV coincide by construction.
At longer horizons the matched selector performs better in aggregate,
particularly under serially correlated disturbances. The results support the benefit of horizon matching within
this controlled design. They do not imply that every scenario
favors the proposed rule, or that the ranking holds under
all possible data-generating mechanisms. No branch tracking
is used in this experiment.

\begin{table}[t]
\centering
\caption{Frozen pooled-selector simulation (Checkpoint 03). Geometric RMSFE
ratios below one favor horizon-matched forecast-CV.}
\label{tab:cp03}
\begin{tabular}{@{}rcc@{}}
\toprule
$h$ & Relative to one-step FCV & Relative to GCV \\
\midrule
1  & 1.000 & 0.912 \\
3  & 0.947 & 0.816 \\
6  & 0.861 & 0.688 \\
12 & 0.762 & 0.554 \\
\bottomrule
\end{tabular}
\end{table}

Three figures visualize distinct features of the primary
forecast-CV experiment. \Cref{fig:cp03_loss_surfaces} illustrates
how the chosen horizon changes the forecast-MSE surface over
normalized smoothness. \Cref{fig:cp03_horizon_noise} compares
horizon-matched selection with one-step tuning across the noise
mechanisms. \Cref{fig:cp03_oracles} distinguishes the feasible
predictive choice from two simulation-only latent oracles.
The latent oracles are diagnostics, not deployable selection
rules.

\begin{figure}[t]
\centering
\includegraphics[width=\columnwidth]{figures/fig_sim05_objective_curves.pdf}
\caption{Checkpoint 03 example of horizon-specific forecast-loss
surfaces over normalized smoothness, with each curve scaled by its
minimum. The figure uses a fixed synthetic scenario with oscillatory
trend and AR(1) observation noise. It illustrates the effect of the
forecast horizon on the objective, not outer-test performance.}
\label{fig:cp03_loss_surfaces}
\end{figure}

\begin{figure}[t]
\centering
\includegraphics[width=\columnwidth]{figures/fig_sim01_horizon_matching_by_noise.pdf}
\caption{Checkpoint 03 geometric forecast-RMSFE ratios of
horizon-matched forecast-CV to one-step forecast-CV,
stratified by observation-noise mechanism. Intervals are
seed-resampling uncertainty summaries from the frozen
simulation; values below one favor horizon matching.}
\label{fig:cp03_horizon_noise}
\end{figure}

\begin{figure}[t]
\centering
\includegraphics[width=\columnwidth]{figures/fig_sim04_smoothness_targets.pdf}
\caption{Checkpoint 03 mean selected smoothness across forecast
horizons for feasible forecast-CV, a latent historical
trend-recovery oracle, and a latent future-forecast oracle.
Both oracles use unobserved simulation information and are
included solely to distinguish statistical targets.}
\label{fig:cp03_oracles}
\end{figure}


\subsection{Distinct recovery and prediction objectives}

The simulated latent-trend oracles displayed in
\cref{fig:cp03_oracles} clarify that historical
signal recovery and extrapolation are
different optimization targets. Both
oracles use unobserved latent information
and are not implementable alternatives to
historical forecast-CV. A smoother can
recover much of the historical trend yet
give unstable terminal differences when
continued into the future.

With difference order $d=2$, all compared
PLS forecasts use linear continuation.
The experiment holds this forecast class
fixed while varying the smoothness-selection
criterion; it does not establish the best
difference order or stochastic residual model.
